\documentclass[UTF8,a4paper,12pt]{ctexart}

\usepackage{geometry}
\usepackage{booktabs}
\usepackage{array}
\usepackage{tabularx}
\usepackage{makecell}
\usepackage{enumitem}
\usepackage{xcolor}
\usepackage{listings}
\usepackage{hyperref}
\usepackage{amsmath}
\usepackage{caption}

\geometry{left=2.6cm,right=2.6cm,top=2.6cm,bottom=2.6cm}

\setlist[itemize]{leftmargin=2em}
\setlist[enumerate]{leftmargin=2em,itemsep=0.25em,topsep=0.4em}

\hypersetup{
    colorlinks=true,
    linkcolor=black,
    urlcolor=blue,
    citecolor=black
}

\definecolor{codebg}{RGB}{248,248,248}
\definecolor{codegray}{RGB}{90,90,90}

\lstdefinestyle{codestyle}{
    backgroundcolor=\color{codebg},
    basicstyle=\ttfamily\small,
    breaklines=true,
    breakatwhitespace=false,
    columns=fullflexible,
    keepspaces=true,
    showstringspaces=false,
    frame=single,
    rulecolor=\color{codegray},
    tabsize=2
}

\lstset{style=codestyle}

\title{\textbf{具身视觉搜索 Agent 设计文档}}
\author{}
\date{}

\begin{document}

\maketitle



\section{任务定义}

本系统面向 AI2-THOR 室内场景中的具身视觉搜索任务。给定一个陌生室内环境和用户自然语言指令，Agent 基于第一人称视觉观察、可交互物体列表、历史动作反馈和搜索记忆，逐步完成目标物体的定位与必要交互。

系统支持两类任务输入。第一类为纯文本指令，例如“找到桌上的杯子”；第二类为图文联合指令，即用户在房间视角中点选目标物体后，系统生成目标截图，并将其作为参考图像输入给 Agent。Agent 的输出为结构化决策结果，包括当前决策摘要、搜索记忆更新和下一步高层动作。

该任务要求模型在受限动作空间内完成高层决策模拟器命令，而是通过多轮“观察—决策—执行—反馈”的交互完成搜索过程。

\section{总体架构}

本系统采用 LLM 驱动的多模态循环决策架构。核心流程如图 \ref{fig:agent-loop} 所示。

\begin{center}
\begin{minipage}{0.78\textwidth}
\centering
\begin{tabular}{c}
环境观察 \\
$\downarrow$ \\
多模态 Prompt 构造 \\
$\downarrow$ \\
OpenAI-compatible 多模态模型调用 \\
$\downarrow$ \\
结构化输出解析 \\
$\downarrow$ \\
高层动作执行 \\
$\downarrow$ \\
反馈、记忆与轨迹更新 \\
$\downarrow$ \\
进入下一轮决策
\end{tabular}
\captionof{figure}{Agent 决策循环}
\label{fig:agent-loop}
\end{minipage}
\end{center}

系统主要由五类模块组成，如表 \ref{tab:modules} 所示。

\begin{center}
\begin{minipage}{0.95\textwidth}
\captionof{table}{系统模块划分}
\label{tab:modules}
\begin{tabularx}{\textwidth}{p{3.1cm}X}
\toprule
模块 & 功能 \\
\midrule
环境与 API 层 & 封装 AI2-THOR 场景加载、机器人视角、房间视角、Agent reset / step / run / stop 等接口。 \\
决策核心层 & 负责组织 observation、调用模型、解析动作、执行动作并更新状态。 \\
Prompt 与模型层 & 将任务、视觉输入、可见物体、记忆和轨迹组织为多模态模型输入。 \\
动作执行层 & 将高层动作映射到 AI2-THOR 可执行动作，并进行可见性与 affordance 检查。 \\
记忆与轨迹层 & 维护搜索记忆、失败反馈、已见物体集合和完整交互轨迹。 \\
\bottomrule
\end{tabularx}
\end{minipage}
\end{center}

其中，\texttt{LLMEmbodiedAgent.step()} 是单步决策的主入口，负责完成从观察获取、Prompt 构造、模型调用、输出解析到动作执行的完整链路。

\section{决策流程}

\subsection{初始化}

任务开始时，系统调用 Agent reset。初始化阶段写入用户指令、最大步数和可选目标参考图像，清空上一轮搜索状态与轨迹记录，并从当前环境获取初始 observation。初始化完成后，Agent 进入可执行状态。

\subsection{单步决策}

每次 step 包含以下过程：

\begin{enumerate}
    \item 获取当前第一人称图像和环境元信息；
    \item 整理当前可见且可交互的物体；
    \item 读取上一轮动作反馈、搜索记忆和最近轨迹；
    \item 构造多模态 Prompt；
    \item 调用 OpenAI-compatible 多模态模型；
    \item 解析模型输出的 JSON；
    \item 执行高层动作；
    \item 根据执行结果更新记忆、状态和轨迹。
\end{enumerate}

\subsection{连续运行}

连续运行模式下，系统循环执行 step，直到出现以下条件之一：Agent 输出 \texttt{end}，达到最大步数，用户手动停止，模型输出无法解析，或动作执行出现严重异常。

\section{观测表示}

Agent 的输入由视觉信息和结构化状态共同构成。

视觉信息包括当前第一人称机器人视角，以及可选的目标参考图像。若上一轮执行了 \texttt{observe}，系统会额外提供由同一位置采集的多方向观察图像。实际输入中，当前机器人图像承担 front 方向，上一轮 observe 生成的 left / back / right 图像作为补充视角输入，从而避免重复发送同一 front 图像。

结构化状态包括表 \ref{tab:observation} 中的字段。

\begin{center}
\begin{minipage}{0.95\textwidth}
\captionof{table}{结构化观测字段}
\label{tab:observation}
\begin{tabularx}{\textwidth}{p{4.2cm}X}
\toprule
字段 & 含义 \\
\midrule
\makecell[l]{\texttt{visible\_}\\\texttt{interactable\_}\\\texttt{objects}} & 当前可见且可交互的物体引用。 \\
\texttt{holding\_objects} & 当前已持有物体。 \\
\texttt{last\_action\_feedback} & 上一步动作执行结果。 \\
\texttt{memory} & 搜索过程中的轻量语义记忆。 \\
\texttt{recent\_steps} & 最近若干步轨迹摘要。 \\
\texttt{target\_reference} & 可选目标参考图像说明。 \\
\bottomrule
\end{tabularx}
\end{minipage}
\end{center}

可见物体不会直接暴露 AI2-THOR 内部 objectId，而是转换为 \texttt{Cup\_1}、\texttt{Drawer\_2} 等可读引用，并附带距离等级和可用 affordance。模型只能使用这些可见引用进行导航或交互，从而降低幻觉动作和非法对象操作的概率。

\section{输出格式}

模型输出采用固定 JSON 结构：

\noindent\textbf{模型输出格式}
\begin{lstlisting}
{
  "thought": {
    "phase": "visual_search",
    "situation_analysis": "...",
    "spatial_reasoning": "...",
    "memory_reasoning": "...",
    "verification": "...",
    "decision": "..."
  },
  "memory_update": {
    "checked": "...",
    "ruled_out": "...",
    "clue": "...",
    "avoid": "..."
  },
  "action": {
    "name": "move forward*3",
    "argument": null,
    "confidence": 0.8
  }
}
\end{lstlisting}

其中，\texttt{thought} 记录简短、可调试的决策摘要，不要求模型输出完整推理链；\texttt{memory\_update} 用于更新搜索状态，包括已检查区域、已排除区域、新线索和应避免重复的行为；\texttt{action} 表示下一步高层动作，由执行层进一步映射到 AI2-THOR 动作。

输出解析器会校验 \texttt{thought} 与 \texttt{action} 字段，并兼容部分轻微格式异常。例如，解析器能够从 Markdown 代码块或混合文本中提取第一个 JSON 对象，也能够对部分截断 JSON 进行修复；同时，解析器会对动作名称、重复次数和旧字段格式进行归一化。如果模型输出仍无法解析，系统会进入解析失败恢复流程。

\section{高层动作空间}

系统将动作空间划分为五类，如表 \ref{tab:actions} 所示。

\begin{center}
\begin{minipage}{0.95\textwidth}
\captionof{table}{高层动作空间}
\label{tab:actions}
\begin{tabularx}{\textwidth}{p{3.2cm}X}
\toprule
类型 & 动作 \\
\midrule
移动与视角调整 & \texttt{move forward}、\texttt{move back}、\texttt{move left}、\texttt{move right}、\texttt{rotate left}、\texttt{rotate right}、\texttt{look up}、\texttt{look down}。 \\
主动观察 & \texttt{observe}。 \\
局部导航 & \texttt{navigate to}。 \\
物体交互 & \texttt{open}、\texttt{close}、\texttt{pickup}、\texttt{put in}、\texttt{toggle}。 \\
结束 & \texttt{end}。 \\
\bottomrule
\end{tabularx}
\end{minipage}
\end{center}

移动和视角动作支持重复表达，例如 \texttt{move forward*3}、\texttt{move right*2}。Prompt 将重复动作主要限制为短连续调整：当方向明确、路径较清楚且只是短距离连续移动时，模型可以使用 \texttt{*2} 或 \texttt{*3}，以避免逐步重复输出同一动作；当路径不确定、靠近障碍物或上一轮同方向移动失败时，模型应避免使用多步重复动作。

导航和交互动作必须使用当前可见物体引用作为参数。执行层会检查目标是否可见，以及是否具备对应 affordance。例如，\texttt{open} 只能作用于 openable 物体，\texttt{pickup} 只能作用于 pickupable 物体，\texttt{put in} 要求当前手中有物体且目标是 receptacle。

\section{搜索记忆}

具身视觉搜索具有长程性和局部观测限制。为减少重复搜索，系统维护轻量语义记忆：

\noindent\textbf{语义记忆字段}
\begin{lstlisting}
summary
checked
ruled_out
avoid
recent_clues
\end{lstlisting}

其中，\texttt{checked} 记录已检查区域或物体，\texttt{ruled\_out} 记录可排除区域，\texttt{avoid} 记录失败动作或无效路线，\texttt{recent\_clues} 保存近期搜索线索。每一步结束后，系统根据模型输出的 \texttt{memory\_update} 和动作反馈更新记忆。如果动作失败，系统会自动加入避免重复该失败动作的提示。

该设计以文本记忆方式为模型提供有限历史状态，在实现复杂度和搜索稳定性之间取得折中。

\section{目标参考图像机制}

系统支持用户通过房间视角点选目标物体。点选后，环境层基于第三方相机视角的实例分割、检测框和物体属性对候选对象排序，并生成目标物体截图。该截图在 Agent reset 时作为 target reference image 输入给模型。

目标参考图像只作为视觉匹配提示，不作为可执行对象引用。模型仍需在第一人称视角中寻找与参考图像匹配的目标，并在目标进入当前可见物体列表后，使用对应 ref 执行导航或交互。

\section{错误恢复机制}

系统包含三类鲁棒性设计。

第一，模型输出恢复。若模型输出不是合法 JSON，系统会将解析错误和上一次无效输出片段重新加入 Prompt，要求模型返回单个合法 JSON 对象。系统最多进行三次解析重试。

第二，动作合法性检查。所有导航和交互动作都必须作用于当前可见 ref；若目标不可见、参数缺失或 affordance 不匹配，执行层返回 illegal action，并将错误反馈写入下一轮 observation。

第三，重复动作失败反馈。对于 \texttt{move forward*3} 这类重复动作，执行层逐步执行底层动作。如果中途失败，系统会返回已完成步数、失败步数和错误原因，例如“执行了 1/3 步，第 2 步失败”。该反馈会进入下一轮 Prompt，帮助模型避免继续沿阻塞方向移动。

\section{轨迹记录}

每一步执行完成后，系统记录完整轨迹，包括任务、场景、模型输出、结构化 thought、动作、执行结果、记忆更新、可见物体、已见物体集合、持有物体、Agent 位姿、机器人视角图像和上一轮动作反馈。

轨迹默认保存到：

\noindent\textbf{轨迹保存路径}
\begin{lstlisting}
data/trajectories/latest_trajectory.json
\end{lstlisting}



\section{局限}

当前系统已经能完成具身视觉搜索功能，但仍存在以下限制：

\begin{enumerate}
    \item 缺少显式空间地图，长程探索主要依赖文本记忆和最近轨迹；
    \item 语义记忆未记录精确拓扑、坐标或物体分布；
    \item 目标完成主要依赖模型自我验证，缺少独立 verifier；
    \item \texttt{navigate to} 主要面向当前可见物体，无法直接处理不可见目标的长程路径规划；
    \item 最终效果受多模态模型视觉识别和空间理解能力影响较大。
\end{enumerate}



\section{总结}

本系统实现了一个面向 AI2-THOR 室内环境的具身视觉搜索 Agent。多模态语言模型作为高层决策核心，以第一人称图像、可见物体、动作反馈、搜索记忆和近期轨迹为输入，输出结构化思考摘要、记忆更新和高层动作，并通过多轮交互逐步完成目标搜索。



\appendix

\section{Prompt 设计}

\subsection{Prompt 总体目标}

Prompt 的设计目标是将多模态模型约束为具身决策器。其核心约束包括：

\begin{enumerate}
    \item 只能基于当前第一人称图像、可见物体 ref、搜索记忆、最近轨迹和动作反馈进行决策；
    \item 不得假设全局地图、隐藏物体、模拟器内部 objectId 或未观察到的信息；
    \item 导航和交互动作必须使用当前 \texttt{visible\_interactable\_objects} 中的 ref；
    \item 若目标不确定，应通过观察、移动、转向、导航或打开相关容器收集证据；
    \item 若目标可见且可达，应执行相应交互；
    \item 在 \texttt{pickup} 或 \texttt{end} 前必须进行目标验证；
    \item 每次只返回一个合法 JSON 对象，不输出 Markdown 或额外解释。
\end{enumerate}

\subsection{System Prompt 结构}

实际系统中的 System Prompt 主要包含以下规则。

\noindent\textbf{System Prompt 摘要}
\begin{lstlisting}
Role:
You are an embodied visual-search agent in an AI2-THOR indoor room.

Task:
Complete the human instruction through observation-action interaction.

Allowed evidence:
- current first-person image
- optional observe multi-view images
- visible_interactable_objects refs
- search memory
- recent trajectory
- last action feedback

Core rules:
- Do not assume global map, hidden objects, object ids, or simulator metadata.
- Use only visible_interactable_objects refs for navigation and interaction.
- Treat target reference image as a visual hint, not as an executable object ref.
- If the target is visible and reachable, act on it.
- If uncertain, gather more evidence.
- Avoid repeated failed actions or already checked places.
- End only after the task is completed or repeated evidence shows it cannot be completed.

Movement:
- Navigation actions support *N when repeated small steps are clearly needed:
  move forward*2/*3, move back*2/*3, move left*2/*3, move right*2/*3.
- Turn actions may also use *N when useful: turn left*2, turn right*2.
- When direction is clear, the path is clear, and the goal is only a short continuous adjustment,
  prefer a single repeated action instead of outputting the same move step-by-step.
- If a move in that direction just failed, or the path is uncertain,
  do not use multi-step movement in that direction.
- After blocked feedback, do not repeat the same action blindly.

Thought policy:
Return a concise reasoning summary, not a long chain-of-thought.
The phase must be one of:
initial_scan, visual_search, navigation, interaction, recovery, completion_check.

Output:
Return exactly one JSON object. No markdown.
\end{lstlisting}

\subsection{User Prompt Payload}

每一步的 User Prompt 使用 JSON payload 组织，主要字段如下。

\noindent\textbf{User Prompt Payload 摘要}
\begin{lstlisting}
{
  "task": "user instruction or target-reference instruction",
  "actions": ["predefined high-level actions"],
  "movement_policy": {
    "repeatable_actions": [
      "move forward",
      "move back",
      "move left",
      "move right",
      "rotate left",
      "rotate right",
      "look up",
      "look down"
    ],
    "default": "Prefer *2 to *3 for clear repeated navigation or turn adjustments instead of repeating the same action line by line.",
    "cautious": "Use a single step when a direction just failed, the path is uncertain, the robot is near obstacles, or re-observation is needed.",
    "examples": [
      "move forward*3",
      "move back*2",
      "move left*2",
      "move right*3",
      "turn left*2",
      "turn right*2"
    ]
  },
  "observation": {
    "image": "attached first-person robot image",
    "visible_interactable_objects": [],
    "holding": [],
    "last_feedback": {
      "success": true,
      "message": "..."
    },
    "observe_views": {
      "note": "...",
      "views": []
    }
  },
  "memory": {
    "summary": "...",
    "checked": [],
    "ruled_out": [],
    "avoid": [],
    "recent_clues": []
  },
  "search_policy": {
    "initial_scan_observe": "...",
    "visual_absence": "...",
    "container_search": "...",
    "anti_loop": "...",
    "revisit_rule": "...",
    "observe_views": "..."
  },
  "recent_steps": [],
  "examples": [],
  "output": {
    "thought": {},
    "memory_update": {},
    "action": {}
  }
}
\end{lstlisting}

如果用户提供目标参考图像，payload 会增加：

\noindent\textbf{目标参考图像字段}
\begin{lstlisting}
{
  "target_reference": {
    "image": "attached target reference image",
    "type_hint": "...",
    "note": "The user selected this object from the room view. Use it as the visual target reference. It may not be visible in the current first-person image."
  }
}
\end{lstlisting}

\subsection{输出 JSON 约束}

模型必须返回如下结构。

\noindent\textbf{输出 JSON 约束}
\begin{lstlisting}
{
  "thought": {
    "phase": "one of: initial_scan, visual_search, navigation, interaction, recovery, completion_check",
    "situation_analysis": "required concise analysis of current visual and task evidence",
    "spatial_reasoning": "movement or viewpoint reasoning, or null",
    "memory_reasoning": "how memory affects this decision, or null",
    "verification": "target or task verification before pickup or end, or null",
    "decision": "required concise reason for the selected action"
  },
  "memory_update": {
    "checked": "what was checked this step, or null",
    "ruled_out": "place, object, view, or direction that likely does not contain the target, or null",
    "clue": "useful clue found this step; mention left/right/front/back when observe views provide directional evidence, or null",
    "avoid": "what to avoid repeating next, or null"
  },
  "action": {
    "name": "one action from actions; when direction and path are clear, prefer repeated navigation actions such as move forward*3, move right*3, move left*2, move back*2, or turn left*2 instead of repeating one-step moves; avoid multi-step repeats after a failed move or when uncertain",
    "argument": "visible object ref when required, else null",
    "confidence": "number in [0, 1] or null"
  }
}
\end{lstlisting}

\subsection{动作空间约束}

Prompt 中的动作约束可概括为：

\noindent\textbf{Prompt 中的动作空间}
\begin{lstlisting}
Movement:
- move forward
- move back
- move left
- move right
- rotate left
- rotate right
- look up
- look down

Observation:
- observe

Navigation:
- navigate to <visible_ref>

Interaction:
- open <visible_ref>
- close <visible_ref>
- pickup <visible_ref>
- put in <visible_ref>
- toggle <visible_ref>

Termination:
- end
\end{lstlisting}

其中，移动与视角动作允许 \texttt{action*N} 表达。新版 Prompt 更强调短重复动作，即优先使用 \texttt{*2} 到 \texttt{*3}，而不是过长的连续移动。若上一轮同方向移动失败、路径不确定或需要重新观察，则应退回单步动作。

\subsection{observe 多视角说明}

\texttt{observe} 用于在当前位置获取方向性视觉证据。执行后，系统依次记录 front、left、back、right 四个方向的图像，并在结束时回到原始朝向。

进入下一轮 Prompt 时，当前第一人称图像作为 front 方向单独发送；上一轮 observe 的 left / back / right 图像作为补充视角输入。Prompt 要求模型根据这些视角判断下一步应转向、前进还是排除某个方向。

\noindent\textbf{observe 多视角说明}
\begin{lstlisting}
The attached environment images are ordered as:
current robot view/front, observe-left, observe-back, observe-right.

The current robot view is the front direction after observe returned to the original heading.
The observe-left, observe-back, and observe-right images were captured from the same position
after rotating left 90, 180, and 270 degrees.

Use these images to decide which direction to rotate or move next.
Objects seen only in observe-left, observe-back, or observe-right are directional visual clues;
rotate toward that direction before interacting with them using visible refs.

If a target-like object appears in one of these views, record that direction in memory_update.clue,
for example "target-like object appears in left view".
\end{lstlisting}

\subsection{解析失败恢复 Prompt}

当模型输出无法解析时，系统追加恢复提示，核心内容如下。

\noindent\textbf{解析失败恢复 Prompt}
\begin{lstlisting}
Previous output could not be parsed.
Return exactly one valid JSON object:
{
  "thought": {
    "phase": "visual_search",
    "situation_analysis": "...",
    "spatial_reasoning": null,
    "memory_reasoning": null,
    "verification": null,
    "decision": "..."
  },
  "memory_update": {
    "checked": null,
    "ruled_out": null,
    "clue": null,
    "avoid": null
  },
  "action": {
    "name": "observe",
    "argument": null,
    "confidence": 0.0
  }
}
Do not output markdown fences, explanations, or partial JSON.
If a field value is unknown, use null.
\end{lstlisting}

该机制使模型输出错误不会立即导致任务终止，而是优先尝试将模型拉回可执行格式。



\end{document}